% ============================================================
%  tables.tex — Formatted LaTeX tables from analysis results
%  Auto-generated from CSV outputs in results/comparison/
% ============================================================

% ────────────────────────────────────────────────────────────
% TABLE 1: Three-Study Overview
% ────────────────────────────────────────────────────────────
\begin{table}[htbp]
\centering
\caption{Overview of the Three Study Datasets}
\label{tab:three_study_overview}
\small
\begin{tabular}{lllll}
\toprule
 & & \textbf{Study 1} & \textbf{Study 2} & \textbf{Study 3} \\
\midrule
Source        & & StudentLife      & NetHealth           & GLOBEM            \\
University    & & Dartmouth        & Notre Dame          & U. Washington     \\
Year          & & 2013             & 2015--2019          & 2018--2021        \\
$N$           & & 28               & 722                 & 809               \\
BFI measure   & & BFI-44           & BFI-44              & BFI-10            \\
Sensing       & & Smartphone       & Fitbit +            & Smartphone        \\
              & &  (13 modalities) &  Communication      &  (RAPIDS)         \\
Duration      & & 10 weeks         & 4 years             & 1 year (4 cohorts)\\
Depression    & & PHQ-9            & CES-D               & BDI-II, PHQ-4,    \\
              & &                  &                     &  PHQ-9            \\
Anxiety       & & PSS              & STAI, BAI           & GAD-7, PHQ-4      \\
GPA           & & Yes              & Yes                 & No                \\
\bottomrule
\end{tabular}
\begin{tablenotes}
\small
\item \textit{Note.} BFI = Big Five Inventory. RAPIDS = Reproducible Analysis Pipeline for
  Data Streams. PSS = Perceived Stress Scale. STAI = State-Trait Anxiety Inventory.
  BAI = Beck Anxiety Inventory. CES-D = Center for Epidemiologic Studies Depression Scale.
  BDI-II = Beck Depression Inventory-II. PHQ = Patient Health Questionnaire.
  GAD-7 = Generalized Anxiety Disorder 7-item scale.
\end{tablenotes}
\end{table}

% ────────────────────────────────────────────────────────────
% TABLE 2: Grand Synthesis — Personality vs Sensing R² [95% CI]
% Source: results/comparison/supplementary/grand_synthesis.csv
% ────────────────────────────────────────────────────────────
\begin{table}[htbp]
\centering
\caption{Grand Synthesis: Personality vs.\ Sensing $R^2$ with 95\% Confidence Intervals}
\label{tab:grand_synthesis}
\footnotesize
\begin{tabular}{llrllll}
\toprule
\textbf{Study} & \textbf{Outcome} & \textbf{N} &
  \textbf{Personality $R^2$ [95\% CI]} &
  \textbf{Sensing $R^2$ [95\% CI]} &
  \textbf{Winner} & \textbf{$\Delta R^2$} \\
\midrule
% Study 1 (S1) — StudentLife
S1 & PHQ-9        & 27  & $-.55$ [$-2.34$, $.24$]  & $.20$ [$-.65$, $.82$]   & Sensing  & $.75$ \\
S1 & PSS          & 27  & \textbf{$.39$ [$.01$, $.72$]}  & $-.23$ [$-1.12$, $.22$] & \textbf{Pers.} & $.62$ \\
S1 & Loneliness   & 27  & $.20$ [$-.19$, $.50$]    & $-.25$ [$-.69$, $.25$]  & Pers.    & $.46$ \\
S1 & GPA          & 27  & $-.17$ [$-.84$, $.26$]   & $-1.85$ [$-5.58$, $-.28$] & Pers.  & $1.68$ \\
\midrule
% Study 2 (S2) — NetHealth
S2 & CES-D        & 498 & \textbf{$.279$ [$.088$, $.391$]} & $-.011$ [$-.085$, $.051$] & \textbf{Pers.} & $.290$ \\
S2 & STAI-Trait   & 498 & \textbf{$.522$ [$.409$, $.620$]} & $.002$ [$-.090$, $.074$]  & \textbf{Pers.} & $.520$ \\
S2 & BAI          & 498 & \textbf{$.176$ [$-.025$, $.334$]} & $-.035$ [$-.127$, $.017$] & \textbf{Pers.} & $.211$ \\
S2 & Loneliness   & 498 & $.011$ [$-.103$, $.063$] & $-.007$ [$-.089$, $.055$] & Pers.   & $.018$ \\
S2 & Self-Esteem  & 717 & \textbf{$.412$ [$.260$, $.495$]} & $-.008$ [$-.062$, $.046$] & \textbf{Pers.} & $.420$ \\
S2 & GPA          & 220 & $.002$ [$-.320$, $.103$] & $-.044$ [$-.274$, $.050$] & Pers.   & $.046$ \\
\midrule
% Study 3 (S3) — GLOBEM
S3 & BDI-II       & 779 & \textbf{$.091$ [$-.008$, $.167$]} & $-.001$ [$-.089$, $.064$] & \textbf{Pers.} & $.092$ \\
S3 & STAI-State   & 799 & \textbf{$.200$ [$.078$, $.260$]}  & $-.007$ [$-.093$, $.043$] & \textbf{Pers.} & $.208$ \\
S3 & PSS-10       & 800 & \textbf{$.143$ [$-.002$, $.221$]} & $-.009$ [$-.087$, $.032$] & \textbf{Pers.} & $.152$ \\
S3 & CESD         & 800 & \textbf{$.093$ [$-.000$, $.167$]} & $-.010$ [$-.083$, $.032$] & \textbf{Pers.} & $.103$ \\
S3 & UCLA         & 798 & \textbf{$.092$ [$.005$, $.149$]}  & $-.023$ [$-.082$, $.016$] & \textbf{Pers.} & $.115$ \\
\bottomrule
\end{tabular}
\begin{tablenotes}
\small
\item \textit{Note.} $R^2$ values estimated via cross-validated ridge regression; negative $R^2$
  indicates prediction worse than intercept-only model. 95\% CIs obtained by bootstrap
  ($B = 1{,}000$). Bold indicates personality $R^2$ CI excludes zero.
  Pers.\ = Personality; $\Delta R^2$ = Personality $R^2$ $-$ Sensing $R^2$.
  Personality wins in 14 of 15 comparisons.
\end{tablenotes}
\end{table}

% ────────────────────────────────────────────────────────────
% TABLE 3: Meta-Analysis
% Source: results/comparison/meta_analysis.csv
% ────────────────────────────────────────────────────────────
\begin{table}[htbp]
\centering
\caption{Random-Effects Meta-Analysis: Neuroticism--Mental Health Correlations Across Studies}
\label{tab:meta_analysis}
\small
\begin{tabular}{lrrlrr}
\toprule
\textbf{Trait--Outcome} & \textbf{$k$} & \textbf{$N$} &
  \textbf{Pooled $r$ [95\% CI]} & \textbf{$I^2$ (\%)} & \textbf{$p$} \\
\midrule
N $\to$ Depression       & 3 & 1{,}304 & \textbf{$.444$ [$.235$, $.614$]} & 91.8 & $< .001$ \\
N $\to$ Anxiety/Stress   & 3 & 1{,}324 & \textbf{$.632$ [$.384$, $.795$]} & 96.2 & $< .001$ \\
N $\to$ Loneliness       & 3 & 1{,}323 & $.135$ [$-.147$, $.396$]         & 94.4 & $.348$   \\
E $\to$ Loneliness       & 3 & 1{,}323 & $-.144$ [$-.465$, $.209$]        & 96.6 & $.425$   \\
C $\to$ Depression       & 3 & 1{,}304 & \textbf{$-.209$ [$-.291$, $-.124$]} & 44.8 & $< .001$ \\
N $\to$ Perceived Stress & 2 &    827  & \textbf{$.576$ [$.071$, $.846$]} & 88.1 & $.028$   \\
C $\to$ GPA              & 2 &    248  & \textbf{$.376$ [$.064$, $.620$]} & 64.0 & $.019$   \\
\bottomrule
\end{tabular}
\begin{tablenotes}
\small
\item \textit{Note.} $k$ = number of studies; $I^2$ = heterogeneity index (proportion of
  variance due to between-study differences); N = Neuroticism; E = Extraversion;
  C = Conscientiousness. Bold indicates $p < .05$.
  High $I^2$ values reflect genuine between-site variation rather than methodological
  inconsistency. Random-effects model (DerSimonian--Laird).
\end{tablenotes}
\end{table}

% ────────────────────────────────────────────────────────────
% TABLE 4: Clinical Classification
% Source: results/comparison/clinical_classification.csv
% Key comparisons: Pers-only vs Pers+Beh vs Beh-only
% ────────────────────────────────────────────────────────────
\begin{table}[htbp]
\centering
\caption{Clinical Classification Performance: Personality vs.\ Sensing Feature Sets}
\label{tab:clinical_classification}
\footnotesize
\begin{tabular}{llllll}
\toprule
\textbf{Study / Outcome} & \textbf{Features} &
  \textbf{AUC [95\% CI]} &
  \textbf{Sensitivity [95\% CI]} &
  \textbf{Specificity [95\% CI]} \\
\midrule
\multicolumn{5}{l}{\textit{Study 2 — CES-D $\geq 16$}} \\
 & Pers-only   & \textbf{$.751$ [$.622$, $.882$]} & $.849$ [$.500$, $1.00$] & $.634$ [$.364$, $.861$] \\
 & Pers+Beh    & \textbf{$.833$ [$.698$, $.932$]} & $.839$ [$.495$, $1.00$] & $.770$ [$.489$, $1.00$] \\
 & Beh-only    & $.552$ [$.402$, $.704$]          & $.648$ [$.139$, $1.00$] & $.607$ [$.112$, $.982$] \\
\midrule
\multicolumn{5}{l}{\textit{Study 2 — STAI-Trait $\geq 45$}} \\
 & Pers-only   & \textbf{$.795$ [$.676$, $.896$]} & $.855$ [$.615$, $1.00$] & $.705$ [$.487$, $.898$] \\
 & Pers+Beh    & \textbf{$.856$ [$.707$, $.960$]} & $.852$ [$.556$, $1.00$] & $.798$ [$.571$, $1.00$] \\
 & Beh-only    & $.577$ [$.416$, $.749$]          & $.702$ [$.200$, $1.00$] & $.592$ [$.160$, $.963$] \\
\midrule
\multicolumn{5}{l}{\textit{Study 3 — BDI-II $\geq 14$}} \\
 & Pers-only   & \textbf{$.654$ [$.552$, $.763$]} & $.708$ [$.323$, $.985$] & $.604$ [$.265$, $.919$] \\
 & Pers+Beh    & \textbf{$.671$ [$.520$, $.778$]} & $.732$ [$.373$, $.938$] & $.632$ [$.374$, $.890$] \\
 & Beh-only    & $.569$ [$.446$, $.658$]          & $.680$ [$.228$, $1.00$] & $.542$ [$.124$, $.917$] \\
\midrule
\multicolumn{5}{l}{\textit{Study 3 — BDI-II $\geq 20$}} \\
 & Pers-only   & \textbf{$.667$ [$.548$, $.788$]} & $.718$ [$.368$, $1.00$] & $.623$ [$.361$, $.935$] \\
 & Pers+Beh    & \textbf{$.686$ [$.515$, $.825$]} & $.710$ [$.259$, $.929$] & $.682$ [$.340$, $.957$] \\
 & Beh-only    & $.604$ [$.458$, $.722$]          & $.650$ [$.308$, $1.00$] & $.637$ [$.182$, $.957$] \\
\midrule
\multicolumn{5}{l}{\textit{Study 3 — PSS $\geq 20$}} \\
 & Pers-only   & \textbf{$.686$ [$.575$, $.773$]} & $.721$ [$.404$, $.937$] & $.633$ [$.366$, $.914$] \\
 & Pers+Beh    & \textbf{$.704$ [$.594$, $.810$]} & $.680$ [$.433$, $.935$] & $.718$ [$.383$, $.931$] \\
 & Beh-only    & $.605$ [$.484$, $.712$]          & $.664$ [$.233$, $.966$] & $.590$ [$.262$, $.933$] \\
\bottomrule
\end{tabular}
\begin{tablenotes}
\small
\item \textit{Note.} AUC = area under the receiver operating characteristic curve.
  Pers-only = Big Five personality traits only; Beh-only = passive sensing behavioral
  features only; Pers+Beh = combined model. CIs obtained by bootstrap ($B = 1{,}000$).
  Best model selected by cross-validated AUC (LR = logistic regression, RF = random forest).
  Bold indicates AUC $> .65$.
\end{tablenotes}
\end{table}

% ────────────────────────────────────────────────────────────
% TABLE 5: Incremental Validity
% Source: results/comparison/incremental_validity.csv
% ────────────────────────────────────────────────────────────
\begin{table}[htbp]
\centering
\caption{Incremental Validity: Does Sensing Add Predictive Value Beyond Personality?}
\label{tab:incremental_validity}
\small
\begin{tabular}{llrrrrrrl}
\toprule
\textbf{Study} & \textbf{Outcome} & \textbf{N} &
  \textbf{$R^2_{\mathrm{pers}}$} &
  \textbf{$R^2_{\mathrm{full}}$} &
  \textbf{$\Delta R^2$} &
  \textbf{$F$} & \textbf{$p$} & \textbf{Sig.\textsuperscript{a}} \\
\midrule
S2 & CES-D     & 363 & .357 & .367 & .009 & 1.75 & $.157$ & \\
S2 & STAI-Trait & 363 & .568 & .570 & .002 & 0.68 & $.566$ & \\
S2 & BAI       & 363 & .245 & .248 & .003 & 0.43 & $.734$ & \\
S3 & BDI-II    & 588 & .130 & .154 & .024 & 3.31 & $.006$ & \textbf{*} \\
S3 & STAI-State & 593 & .226 & .244 & .018 & 2.71 & $.020$ & \\
S3 & PSS-10    & 594 & .176 & .190 & .014 & 1.99 & $.078$ & \\
S3 & CESD      & 594 & .112 & .130 & .018 & 2.39 & $.036$ & \\
S3 & UCLA Loneliness & 592 & .108 & .116 & .008 & 1.08 & $.370$ & \\
\bottomrule
\end{tabular}
\begin{tablenotes}
\small
\item \textit{Note.} $R^2_{\mathrm{pers}}$ = $R^2$ from personality-only model;
  $R^2_{\mathrm{full}}$ = $R^2$ from personality + sensing model;
  $\Delta R^2$ = incremental variance explained by sensing features.
  $F$-test evaluates whether sensing features add significant predictive value over
  personality alone. \textsuperscript{a}Significance after FDR correction ($q < .05$):
  \textbf{*} $p_{\mathrm{FDR}} < .05$.
  S2 has 3 behavioral features; S3 has 5 behavioral features.
\end{tablenotes}
\end{table}
